\chapter{Protocollo e metriche}\label{ch:metodo}
\begingroup
\setlength{\parskip}{0pt}
\setlength{\abovedisplayskip}{5pt plus 1pt}
\setlength{\belowdisplayskip}{5pt plus 1pt}
\setlength{\abovedisplayshortskip}{3pt plus 1pt}
\setlength{\belowdisplayshortskip}{3pt plus 1pt}
\section{Split cronologici ed eleggibilità}
La Tabella~\ref{tab:split} riporta gli split ufficiali, con fine inclusiva all'ultimo microsecondo del giorno. Il primo evento di sessione è il decision point.
\begin{table}[htbp]\centering
\begin{tabular}{lll}\toprule Fase & Inizio & Fine\\\midrule
Warm-up & 30/10/2005 & 31/12/2008\\
Validation & 01/01/2009 & 28/02/2009\\
Test held-out & 01/03/2009 & 05/05/2009\\\bottomrule
\end{tabular}
\caption{Intervalli cronologici congelati. Per valutare una decisione deve essere completa anche la finestra futura di sette giorni.}\label{tab:split}
\end{table}

La separazione cronologica riproduce la disponibilità progressiva dell'informazione: il warm-up costruisce storia e scala delle feature, la validation sceglie i parametri, il test valuta le configurazioni congelate su un periodo successivo. Uno split casuale non riprodurrebbe questo ordine e potrebbe far dipendere la preparazione del modello da eventi posteriori alle decisioni valutate. È questa esigenza sequenziale, non una presunta stazionarietà dei gusti, a motivare il protocollo.

Il warm-up delimita gli eventi storici e i decision point dei contesti utilizzati per il fitting dello scaler. Lo stato non azzera la storia al 30/10/2005: eventuali eventi precedenti contribuiscono a $D_{<t}$, ma un decision point esterno al warm-up non genera righe di fitting.

Per una fase $F=[s_F,e_F]$, il numero di decisioni utilizzabili è
\begin{equation}
N_{F,u}^{\mathrm{usable}}=
\#\{t\in\mathcal S_u\cap F:t+\Delta\leq e_F,\
|\mathcal A_u(t)|>0\}.\label{eq:usable}
\end{equation}
Il runner richiede che il pool sia non vuoto a tutte le profondità confrontate; $N_{\mathrm{train},u}$ conta invece \emph{eventi storici} del warm-up, non sessioni. L'utente è eleggibile se tutti e tre i conteggi raggiungono dieci. I risultati ufficiali includono 562 utenti, 77.399 decisioni di validation e 82.905 di test. Un pool vuoto non produce una raccomandazione e non contribuisce alle metriche. La fine del periodo non viene estesa per recuperare decisioni prive di futuro completo.

La coorte comune è selezionata retrospettivamente: l'attività del test ne determina l'appartenenza, influenzando anche la popolazione confrontata in validation. Ciò non immette ascolti futuri nelle feature né metriche held-out nel tuning; condiziona però i risultati a utenti attivi in tutti i periodi, senza generalizzazione automatica a utenti occasionali o inattivi. Poiché le sessioni sono realmente osservate, si valuta condizionatamente al ritorno dell'utente, senza misurare se la raccomandazione lo induca a tornare.

La soglia di dieci è un vincolo operativo per confrontare profili con più osservazioni utilizzabili; non è una necessità matematica del reward né una soglia di affidabilità statistica dimostrata.

\section{Validation, congelamento e test}
Le cinque famiglie generano 12 configurazioni per $N$: tre baseline senza sweep, tre valori di $\varepsilon$ e sei di $\alpha$. Le griglie sono $\varepsilon\in\{0{,}05;0{,}10;0{,}20\}$ e $\alpha\in\{0;0{,}1;0{,}25;0{,}5;1;2\}$. Moltiplicando per $N=50,100,200$ si ottengono 36 run.

$N=100$ è lo scenario principale definito e documentato prima dell'esecuzione della validation. Il commit \metric{c716731} del 7 ottobre 2026 lo fissa nel README del runner e nel selettore, prima dell'avvio ufficiale del 19 ottobre (Appendice~\ref{app:repro}). $N=50/200$ restano sensibilità. Per ogni famiglia si massimizza NDCG@10 a $N=100$; a parità si preferiscono reward maggiore e parametro minore. Il manifesto congela $\varepsilon=0{,}05$, $\alpha=2$, seed 42 per Random ed $\varepsilon$-Greedy, $K=(5,10,20)$. Il test non seleziona $N$ o iperparametri.

\par\noindent\begin{minipage}{\linewidth}
Come giustificazione tecnica, $N=100$ è intermedio fra gli scenari confrontati e limita l'unione a 300 item, contro 150 e 600 agli estremi. Bilancia ampiezza ammissibile e lavoro per decisione; non esiste nelle evidenze esaminate un benchmark preliminare che lo dimostri ottimale. La scelta resta operativa e antecedente ai risultati.
\end{minipage}\par

NDCG@10 fornisce un criterio univoco di selezione sensibile alla posizione degli artisti pertinenti, evitando di scegliere a posteriori la metrica più favorevole. La domanda scientifica rimane multidimensionale: reward Top-1, novelty, ILD e coverage descrivono aspetti che NDCG non riassume. Non sono quindi marginali all'obiettivo del progetto, pur non guidando il tuning salvo lo spareggio previsto sulla reward. I cutoff 5 e 20 mostrano se i comportamenti cambiano accorciando o allungando lo stesso ranking; 10 resta il riferimento principale. Il test verifica quanto emerge dopo il congelamento, senza riaprire tale scelta.

Il warm-up aggiorna solo storia e FeatureScaler: nessuna azione, reward, coda di feedback o apprendimento di LinUCB. All'inizio di ogni fase \metric{run_phase()} resetta la policy: $A=I_{15}$, $b=0_{15}$ per LinUCB e seed ripristinato per le policy casuali. LinUCB apprende poi dai soli feedback maturi della fase valutata. Le baseline a conteggi non apprendono dai reward.

Nel test, istanze indipendenti ricostruiscono anche gli ascolti di validation anteriori alla prima decisione, senza importarvi matrici o feedback della validation. La storia musicale è quindi distinta dalla memoria di apprendimento. La coda deve essere vuota all'avvio della fase; il vincolo $t+7d\leq e_F$ permette di maturare tutti i feedback entro la fine, quando vengono applicati quelli residui senza nuove scelte. Lo scaler resta congelato dal solo warm-up, con digest verificato.

\section{Spazio delle raccomandazioni e rilevanza}
Sia $L^K_{u,t}$ il prefisso di massimo $K$ \emph{tracce} e sia $G^K_{u,t}=(g_1,\ldots,g_m)$ la sequenza di artisti ottenuta preservando la prima occorrenza. Vale $m\leq K$. Non si cerca una nuova Top-K artisti nel pool. Il riferimento futuro è $R_{u,t}=\{g(e):u(e)=u,\ t<\tau(e)\leq t+\Delta\}$.
La rilevanza è binaria, $\mathrm{rel}_j=\mathbb1[g_j\in R_{u,t}]$. Precision e Recall sono
\begin{equation}
\mathrm{Precision@K}=\frac{\sum_{j=1}^m\mathrm{rel}_j}{m},
\qquad
\mathrm{Recall@K}=\frac{|G^K_{u,t}\cap R_{u,t}|}{|R_{u,t}|}.
\end{equation}
Se il riferimento futuro è vuoto, Recall è indefinita ed esclusa dall'aggregazione; la precisione di una lista non vuota è zero. NDCG segue la normalizzazione del guadagno scontato \cite{jarvelin2002}:
\begin{align}
\mathrm{DCG@K}&=\sum_{j=1}^{m}\frac{\mathrm{rel}_j}{\log_2(j+1)},\\
\mathrm{IDCG@K}&=\sum_{j=1}^{\min(K,|R_{u,t}|)}\frac1{\log_2(j+1)},
&\mathrm{NDCG@K}&=\frac{\mathrm{DCG@K}}{\mathrm{IDCG@K}}.\label{eq:ndcg}
\end{align}
NDCG è esclusa quando IDCG è zero. Il denominatore usa $K$ e il numero di artisti futuri rilevanti, anche se i duplicati riducono la lista raccomandata a $m<K$. Non usa solo il numero di artisti rappresentati nella lista e non limita il riferimento agli artisti candidabili. La deduplicazione impedisce che DCG superi IDCG; il codice applica inoltre un limite numerico superiore pari a uno.

Per esempio, con $K=5$, le tracce $(A_1,A_2,B_1,C_1,B_2)$ diventano gli artisti $(A,B,C)$, con $m=3$. Se $R_{u,t}=\{A,C,D\}$, Precision@5 e Recall@5 valgono entrambe $2/3$. DCG vale $1+1/\log_2 4=1{,}5$, IDCG vale $1+1/\log_2 3+1/\log_2 4=2{,}130930$ e NDCG@5 è $0{,}703918$. La normalizzazione dipende dunque anche dalle ripetizioni d'artista e dal riferimento futuro; composizione del pool e capacità di ordinamento concorrono al risultato.

\clearpage
\section{Novelty, diversità, reward e coverage}
Novelty e diversità sono dimensioni distinte dall'accuratezza, come discusso in letteratura \cite{vargas2011}; qui si adottano convenzioni operative specifiche, non tutte le metriche di quel lavoro. La track novelty media separatamente ogni item raccomandato, mentre la artist novelty opera sulla lista deduplicata:
\begin{equation}
\mathrm{TN@K}=\frac1{|L^K|}\sum_{i\in L^K}\mathbb1[C_{u,i}(t)=0],
\qquad
\mathrm{AN@K}=\frac1m\sum_{g\in G^K}\mathbb1[C_{u,g}(t)=0].
\end{equation}
I flag sono salvati prima della decisione, non ricostruiti dalla storia futura. La ILD considera invece tutte le posizioni delle Top-K tracce e una distanza binaria fra artisti. Con $k'=|L^K|$ e $n_g$ numero di tracce dell'artista $g$,
\begin{equation}
\mathrm{ILD@K}=1-\frac{\sum_g n_g(n_g-1)}{k'(k'-1)}\quad(k'\geq2),
\end{equation}
ed è zero per liste più brevi. La valutazione della lista nel suo insieme è centrale anche nei lavori sulla diversificazione \cite{ziegler2005}; la nostra distanza binaria rappresenta soltanto varietà d'identità, non similarità tematica, acustica o semantica, e non applica un reranking di diversificazione.

Average Reward aggrega la reward della sola Top-1 dell'Equazione~\ref{eq:reward}. Non è la media di K reward e non dipende dal cutoff con cui viene osservata la lista. Per le metriche di decisione, la procedura ufficiale prima calcola la media sulle decisioni valide di ciascun utente, poi media fra gli utenti:
\begin{equation}
\overline M=\frac1{|\mathcal U_M|}\sum_{u\in\mathcal U_M}
\frac1{|T_{u,M}|}\sum_{t\in T_{u,M}}M_{u,t}.\label{eq:macro}
\end{equation}
I sottoinsiemi validi dipendono dalla metrica; non si imputano a zero quantità indefinite.

La coverage è globale di fase. Indicando con $I_F$ l'unione degli item candidati nelle decisioni valutabili e con $G_F$ i relativi artisti,
\begin{equation}
C_I@K=\frac{|\bigcup_{u,t}L^K_{u,t}|}{|I_F|},\qquad
C_G@K=\frac{|\bigcup_{u,t}G^K_{u,t}|}{|G_F|}.
\end{equation}
\endgroup
